\subsection{Problema dei due corpi}
\[
\mathcal{H} = \frac{{\vec{p}_1}^2}{2m_1} + \frac{{\vec{p}_2}^2}{2m_2} + V(\vec{x_1} - \vec{x_2}), \quad \text{corpi 1 e 2.}
\]
Voglio ridurmi al problema ad un solo corpo e poi passo al potenziale centrale, dove uso le coordinate sferiche.

So che
\[
[\hat{p}_a^{(i)}, \hat{x}_b^{(j)}] = -i\hbar \delta^{ij} \delta^{ab},
\]
\[
[\hat{p}_a^{(i)}, \hat{p}_b^{(j)}] = [\hat{x}_a^{(i)}, \hat{x}_b^{(j)}] = 0.
\]
Poi introduco
\[
\begin{dcases}
\vec{r} = \vec{x}_1 - \vec{x}_2 \\
\vec{R} = \frac{m_1 \vec{x}_1 + m_2 \vec{x}_2}{m_1 + m_2}
\end{dcases} \quad \text{come nuove coordinate}
\]
e poi costruisco i momenti coniugati tali che
\[
\begin{cases}
[\hat{p}^{(i)}, \hat{r}^{(j)}] = -i\hbar\delta_{ij} = [P^{(i)}, R^{(j)}] \\
[\hat{p}^{(i)}, R^{(j)}] = [P^{(i)}, r^{(j)}] = 0
\end{cases}.
\]
La scelta che funziona è
\[
\begin{dcases}
\hat{\vec{p}} = \frac{m_2 \vec{p}_1 - m_1 \vec{p}_2}{m_1 + m_2} \\
\vec{P} = \hat{\vec{p}}_1 + \hat{\vec{p}}_2
\end{dcases}
\]
che verificano le regole di commutazione (cambio di variabili canonico). In generale
\[
\begin{cases}
\vec{x_1'} = \vec{r} \\
\vec{x_2'} = \vec{R}
\end{cases}, \quad \begin{cases}
\vec{p_1'} = \vec{p} \\
\vec{p_2'} = \vec{P}
\end{cases}
\]
le scrivo come 
\[
\begin{pmatrix}
\vec{x_1'} \\
\vec{x_2'}
\end{pmatrix} = M \begin{pmatrix}
\vec{x_1} \\
\vec{x_2}
\end{pmatrix}, \quad \begin{pmatrix}
\vec{p_1'} & \vec{p_2'}
\end{pmatrix} = N \begin{pmatrix}
\vec{p_1} & \vec{p_2}
\end{pmatrix}, \quad \text{con } N = M^{-1}
\]
\[
\boxed{N = M^{-1} \Leftrightarrow [p_a^{(i)}, p_b^{(j)}] = -i\hbar\delta_{ij}\delta_{ab}} \quad \text{(voglio verificarlo)}.
\]
\begin{align*}
[{p'}_{a'}^{(i)}, {x'}_{b'}^{(j)}] &= [p_a^{(i)}N_{aa'}, M_{b'b} x_b^{(j)}] = \\
&= [p_a^{(i)}, x_b^{(j)}] M_{b'b} N_{aa'} = \\
&= -i\hbar\delta_{ij}\delta_{ab}N_{aa'}M_{b'b} = \\
&= -i\hbar\delta_{ij}M_{b'a}N_{aa'} = \\
&= i\hbar\delta_{ij} (MN)_{b'a'} \stackrel{\text{per ipotesi}}{=} -i\hbar\delta_{ij}(\mathbbm{1})_{b'a'} \Leftrightarrow M = N^{-1}.
\end{align*}
Quindi, come avevo prima,
\[
{p'}_{a'}^{(i)} = -i\hbar\frac{\partial}{\partial {x'}_{a'}^{(i)}},
\]
infatti, per \( (i) \) fissato,
\[
-i\hbar \frac{\partial}{\partial {x'}_{a'}^{(i)}} = -i\hbar \frac{\partial}{\partial x_a^{(i)}} \frac{\partial x_a^{(i)}}{\partial {x'}_{a'}^{(i)}} \quad \Rightarrow \quad M_{aa'} = M_{aa'}^{-1}
\]
\[
\frac{\partial {x'}_{a'}^{(i)}}{\partial x_a^{(i)}} = M_{aa'} \Rightarrow -i\hbar\frac{\partial}{\partial {x'}_{a'}^{(i)}} = -i\hbar \frac{\partial}{\partial x_a^{(i)}} M_{a'a} = p_a^{(i)} M_{aa'}^{-1}
\]
e poi si dimostra che 
\[
\mathcal{H} = \frac{\vec{P}^2}{2M} + \frac{\vec{p}^2}{2\mu} + V(\vec{r}), \quad M = m_1 + m_2, \quad \mu = \frac{m_1 m_2}{m_1 + m_2}.
\]

D'ora in poi studierò solo sistemi ad un corpo, infatti, nel sistema del centro di massa,
\[
\mathcal{H} = \frac{\vec{p}^2}{2m} + V(r)
\]
e tipicamente userò coordinate sferiche 
\[
\begin{cases}
x = r \sin \theta \cos \varphi \\
y = r \sin \theta \sin \varphi \\
z = r \cos \theta
\end{cases}
\]
però adesso devo sepapare \( \vec{p} = -i\hbar \vec{\nabla} \). Per l'impostazione radiale
\[
\hat{p} = \frac{\vec{x}}{\left\Vert \vec{x} \right\Vert} \mathbf{\cdot} \vec{p}
\]
e vedo se vado da qualche parte. Introduco 
\[
\frac{\vec{x}}{\left\Vert \vec{x} \right\Vert} \mathbf{\cdot} \vec{\nabla} = ?
\]
Prendo la componente \( \partial_i \)
\[
\begin{cases}
\partial_i = \frac{\partial}{\partial r} \partial_i r + \frac{\partial}{\partial \theta} \partial_i \theta + \frac{\partial}{\partial \varphi} \partial_i \varphi \\
\frac{\partial \theta}{\partial x_i} x^i = \frac{\partial \varphi}{\partial x_i} x_i = 0
\end{cases} \quad \Rightarrow
\]
\[
\Rightarrow \quad \frac{x_i}{r} \frac{\partial}{\partial r} \frac{\partial}{\partial x_i} |\vec{x}| = \frac{x_i}{r} \frac{\partial}{\partial r} \frac{x_i}{|\vec{x}|} = \frac{x_i^2}{r^2} \frac{\partial}{\partial r} = \frac{\partial}{\partial r}.
\]

Quindi ottengo che \( [\tilde{p}_r, r] = -i\hbar \). Noto che anche \( [\tilde{p}_r + f(r), r] = -i\hbar \). Il secondo termine è critico, quindi 
\[
T = \frac{{\vec{p}}^2}{2m} = \frac{\Delta}{2m} = \sum_i \frac{\partial_i^2}{2m}.
\]
Uso l'equazione classica
\[
(a \mathbf{\cdot} b) = \left\Vert \vec{a} \right\Vert^2 \left\Vert \vec{b} \right\Vert^2 - \left\Vert \vec{a} \wedge \vec{b} \right\Vert^2
\]
e la dimostro, introducendo
\[
\varepsilon^{ijk} = \begin{cases}
0 \quad \text{se due indici hanno lo stesso valore} \\
1 \quad \varepsilon^{123} \\
\text{antisimmetrico rispetto a scambio di indici}
\end{cases} .
\]
Allora
\[
\left( \vec{a} \wedge \vec{b} \right)_i = \varepsilon_{ijk} a_j b_k
\]
\[
\varepsilon^{ijk} \varepsilon^{iab} = \delta^{ia}\delta^{kb} - \delta^{jb} \delta^{ka}
\]
\[
\left\Vert \vec{a} \wedge \vec{b} \right\Vert^2 = \varepsilon^{imn} \varepsilon^{ijk} a_n b_n a_i b_k
\]

Voglio sapere che cos'è \( \hat{\vec{p}}^2 \), partendo dal caso classico, per arrivare a quello quantistico. Uso l'identità vettoriale, che nel caso classico è 
\[
(\vec{a} \mathbf{\cdot} \vec{b})^2 = |\vec{a}|^2 + |\vec{b}|^2 - |\vec{a} \wedge \vec{b}|^2 \quad (\text{da dimostrare}).
\]
Con ciò visto sopra
\[
\varepsilon^{ijk} \varepsilon^{iab} = \delta^{ja} \delta^{kb} - \delta^{jb} \delta^{ka},
\]
\[
(\vec{a} \mathbf{\cdot} \vec{b})_i = \varepsilon_{ijk} a^j b_k
\]
\begin{align*}
|\vec{a} \wedge \vec{b}|^2 &= \varepsilon^{ijk} \varepsilon^{ilm} a^j b^l a^k b^k = \\
&= (\delta^{jl} \delta^{km} - \delta^{jm} \delta^{kl}) a^j b^k a^l b^m = \\
&= a^j a^j b^k b^k - a^j b^j a^k b^k = \\
&= |\vec{a}|^2 + |\vec{b}|^2 - (\vec{a} \mathbf{\cdot} \vec{b})^2
\end{align*}
e ho ciò che cercavo. Poi sostituisco \( \vec{a} \) e \( \vec{b} \) con \( \vec{x} \) e \( \vec{p} \) e ho la relazione classica 
\[
\vec{p}^2 = \vec{p}_r^2 + \frac{\vec{L}^2}{r^2}.
\]

Passo alla quantistica, con \( \hat{\vec{x}} \) e \( \hat{\vec{p}} \) (che non commutano)
\begin{align*}
\varepsilon^{ijk} \varepsilon^{ilm} \hat{x}^j \hat{p}^m &= (\delta^{jl}\delta^{km} - \delta^{jm}\delta^{kl}) \hat{x}^j \hat{p}^k \hat{x}^l \hat{p}^m = \\
&= \hat{x}^j \hat{p}^k \hat{x}^j \hat{p}^k - \hat{x}^j \hat{p}^k \hat{x}^k \hat{p}^j = \\
&= \hat{x}^j \hat{x}^j \hat{p}^k \hat{p}^k + \hat{x}^j [\hat{p}^k, \hat{x}^j] \hat{p}^k - (\hat{x}^j \hat{p}^k \hat{p^j} \hat{x}^k + \hat{x}^j \hat{p}^k [\hat{x}^k, \hat{p}^j]) = \\
&= \hat{r}^2 \hat{\vec{p}}^2 - i\hbar \hat{x}^j \hat{p}^j - (\hat{x}^j \hat{p}^j \hat{p}^k \hat{x}^k + i\hbar \hat{x}^j \hat{p}^j) = \\
&= \hat{r}^2 \hat{\vec{p}}^2 - i\hbar \hat{x}^j \hat{p}^j - ((\hat{\vec{x}} \mathbf{\cdot} \hat{\vec{p}})(\hat{\vec{x}} \mathbf{\cdot} \hat{\vec{p}}) + \hat{x}^j \hat{p}^j [\hat{p}^k, \hat{x}^k] + i\hbar \hat{x}^j \hat{p}^j) = \\
&= \hat{r}^2 \hat{\vec{p}}^2 -i\hbar\hat{\vec{x}} \mathbf{\cdot} \hat{\vec{p}} - (\hat{\vec{x}} \mathbf{\cdot} \hat{\vec{p}}) (\hat{\vec{x}} \mathbf{\cdot} \hat{\vec{p}}) - (-3i\hbar\hat{\vec{x}} \mathbf{\cdot} \hat{\vec{p}} + i\hbar\hat{\vec{x}} \mathbf{\cdot} \hat{\vec{p}}) = \\
&= \hat{r}^2 \hat{\vec{p}}^2 + i\hbar \hat{\vec{x}} \mathbf{\cdot} \hat{\vec{p}} - (\hat{\vec{x}} \mathbf{\cdot} \hat{\vec{p}})^2 = |\vec{x} \wedge \vec{p}|^2 \equiv \vec{L}^2
\end{align*}
dove \( \vec{L} \) ha la proprietà che
\[
\vec{x} \mathbf{\cdot} \vec{L} = \varepsilon^{ijk} x^i x^j p^k = \frac{1}{2} \varepsilon^{ijk} [x^i, x^j] p^k = 0.
\]